% !TEX root = ../main.tex
\corrigesChapitre{4}{Cinétique chimique}
\corrigesSouples

\begin{blocLivre}{cor-4-1}
\begin{corrige}{1}{Lente ou rapide ?}
\demarche{Se demander si l'évolution peut être suivie au cours du temps avec un appareil usuel.}
\begin{questions}
\item Une transformation rapide semble achevée dès le contact des réactifs ; une transformation lente dure de quelques
secondes à plusieurs heures, ce qui permet d'en suivre l'évolution. A et D sont \res{\text{rapides}} ;
B, C et E sont \res{\text{lentes}} (B est même très lente).
\item C : l'absorbance de la solution, car l'ion permanganate est la seule espèce colorée (spectrophotométrie).
E : la pression du gaz dans un récipient fermé, ou le volume de dihydrogène recueilli.
\item Chaque goutte de solution titrante doit réagir aussitôt versée : sinon, la composition du mélange ne
correspondrait pas au volume versé et l'équivalence serait repérée trop tard.
\end{questions}
\end{corrige}
\end{blocLivre}
\begin{blocLivre}{cor-4-2}
\begin{corrige}{2}{Lire un temps de demi-réaction}
\demarche{Lire la valeur finale sur l'asymptote horizontale, puis la date à laquelle la moitié de cette valeur est atteinte.}
\begin{questions}
\item Le temps de demi-réaction est la durée au bout de laquelle l'avancement atteint la moitié de sa valeur finale.
\item $\res{[\ce{I2}]_\mathrm{f}=\qty{4,0}{mmol.L^{-1}}}$ ; $[\ce{I2}]=\qty{2,0}{mmol.L^{-1}}$ pour $\res{t_{1/2}=\qty{15}{min}}$.
\item Le réactif limitant est entièrement consommé et chaque ion peroxodisulfate forme une molécule de diiode :
$\res{[\ce{S2O8^2-}]_0=\qty{4,0}{mmol.L^{-1}}}$.
\item $[\ce{I2}](\qty{30}{min})=\qty{3,0}{mmol.L^{-1}}$, et non $\qty{4,0}{mmol.L^{-1}}$ : la formation du diiode
\res{\text{ralentit}}, car les réactifs se raréfient.
\end{questions}
\schema[corrige]{04-ex02}
\end{corrige}
\end{blocLivre}
\begin{blocLivre}{cor-4-3}
\begin{corrige}{3}{Exploiter des tangentes}
\demarche{Repérer deux points éloignés de chaque tangente, par exemple ses points d'intersection avec les axes.}
\begin{questions}
\item $v_\mathrm{disp}(\mathrm{R})=-\dfrac{\mathrm{d}[\mathrm{R}]}{\mathrm{d}t}$ : c'est l'opposé du coefficient
directeur de la tangente à la courbe à la date considérée.
\item $[\mathrm{R}]=\frac{60}{2}=\qty{30}{mmol.L^{-1}}$ pour $\res{t_{1/2}\approx\qty{21}{min}}$.
\item T$_0$ passe par les points $(0\,;\,60)$ et $(30\,;\,0)$, T$_{1/2}$ par $(21\,;\,30)$ et $(51\,;\,0)$ :
\[v(0)=\frac{60}{30}=\qty{2,0}{mmol.L^{-1}.min^{-1}}=\res{\qty{2,0e-3}{mol.L^{-1}.min^{-1}}}\]
\[v(t_{1/2})=\frac{30}{51-21}=\qty{1,0}{mmol.L^{-1}.min^{-1}}=\res{\qty{1,0e-3}{mol.L^{-1}.min^{-1}}}.\]
\item $\frac{2{,}0}{60}=\frac{1{,}0}{30}\approx\qty{0,033}{min^{-1}}$ : le quotient est \res{\text{le même}}. La vitesse
est proportionnelle à $[\mathrm{R}]$, ce qui est caractéristique d'une loi de vitesse d'ordre 1.
\end{questions}
\schema[corrige]{04-ex03}
\end{corrige}
\end{blocLivre}
\begin{blocLivre}{cor-4-4}
\begin{corrige}{4}{Identifier des facteurs cinétiques}
\demarche{Comparer deux expériences qui ne diffèrent que par un seul paramètre.}
\begin{questions}
\item Concentration : expériences 1 et 2 ; quand $[\ce{I-}]_0$ double, $t_{1/2}$ est divisé par deux : la vitesse
\res{\text{augmente avec la concentration}}. Température : expériences 1 et 3 ; $t_{1/2}$ passe de 12 à 4 min : la
vitesse \res{\text{augmente avec la température}}.
\item Ils accélèrent la réaction sans être consommés : ce sont des \res{\text{catalyseurs}}. Dissous dans la même
phase que les réactifs, ils réalisent une \res{\text{catalyse homogène}}.
\item \res{\text{Oui}} : l'eau oxygénée, réactif limitant, a la même quantité initiale ; ni la température, ni la
concentration du réactif en excès, ni le catalyseur ne modifient l'état final.
\end{questions}
\end{corrige}
\end{blocLivre}
\begin{blocLivre}{cor-4-5}
\begin{corrige}{5}{Un catalyseur solide}
\demarche{Vérifier les propriétés d'un catalyseur : il accélère, il n'est pas consommé, il ne modifie pas l'état final.}
\begin{questions}
\item \ce{2 H2O2(aq) -> 2 H2O(l) + O2(g)}.
\item La réaction, imperceptible sans \ce{MnO2}, devient rapide en sa présence ; la masse du solide est inchangée en fin
de réaction et \ce{MnO2} n'apparaît pas dans l'équation : c'est \res{\text{un catalyseur}}.
\item Le catalyseur est solide, les réactifs sont en solution : \res{\text{catalyse hétérogène}}.
\item La réaction a lieu à la surface du solide : la poudre offre une \res{\text{surface de contact}} bien plus grande
que le grain, donc davantage de sites où la réaction se produit.
\item \res{\text{Non}} : le catalyseur ne modifie pas l'état final ; le volume de dioxygène est fixé par la quantité
de peroxyde d'hydrogène.
\end{questions}
\end{corrige}
\end{blocLivre}
\begin{blocLivre}{cor-4-6}
\begin{corrige}{6}{Suivi par spectrophotométrie}
\demarche{Établir l'état final par un tableau d'avancement, convertir les absorbances en concentrations, puis exploiter la courbe.}
\begin{questions}
\item $n(\ce{H2O2})_0=c_1V_1=\qty{1,0e-4}{mol}$ et $n(\ce{I-})_0=c_2V_2=\qty{4,0e-3}{mol}$, d'où
\[x_{\max}=\min\left(\num{1.0e-4}\,;\,\frac{\num{4.0e-3}}{2}\right)=\qty{1,0e-4}{mol}\ :\quad\res{\text{l'eau oxygénée est limitante}}.\]
$[\ce{I2}]_\mathrm{f}=\dfrac{\num{1.0e-4}}{\num{50.0e-3}}=\res{\qty{2,0e-3}{mol.L^{-1}}}$ et $A_\mathrm{f}=250\times\num{2.0e-3}=\res{0{,}500}$.
\item $[\ce{I2}]=A/k$, soit, en mmol$\cdot$L$^{-1}$ : 0 ; 0,25 ; 0,49 ; 0,69 ; 0,85 ; 1,00 ; 1,30 ; 1,51 ; 1,64 ; 1,76 ; 1,87 ; 1,94.
\item $[\ce{I2}]=\qty{1,0}{mmol.L^{-1}}$ (soit $A=0{,}250$) pour $\res{t_{1/2}=\qty{10}{min}}$.
\item La tangente à l'origine passe par le point $(\qty{14,4}{min}\,;\,\qty{2,0}{mmol.L^{-1}})$, donc
\[v_\mathrm{app}(0)=\frac{\qty{2,0e-3}{mol.L^{-1}}}{\qty{14,4}{min}}\approx\res{\qty{1,4e-4}{mol.L^{-1}.min^{-1}}}.\]
Elle diminue ensuite, car les réactifs se raréfient.
\item D'après les nombres stœchiométriques, $\res{v_\mathrm{disp}(\ce{I-})=2\,v_\mathrm{app}(\ce{I2})}$.
\end{questions}
\schema[corrige]{04-ex06-cor}
\end{corrige}
\end{blocLivre}
\begin{blocLivre}{cor-4-7}
\begin{corrige}{7}{Choisir une méthode de suivi}
\demarche{Pour chaque équation, repérer une espèce colorée, des ions qui apparaissent ou disparaissent, ou un gaz.}
\begin{questions}
\item A : des ions \ce{H3O+} et \ce{Cl-} se forment à partir de molécules : \res{\text{conductimétrie}} (ou pH-métrie).
B : un gaz se forme dans un récipient fermé : \res{\text{mesure de la pression}}.
C : l'ion permanganate est la seule espèce colorée : \res{\text{spectrophotométrie}}.
D : des ions \ce{HO-} sont remplacés par des ions \ce{CH3COO-} : \res{\text{conductimétrie}}.
\item Chaque ion \ce{HO-} ($\lambda=19{,}9$) est remplacé par un ion \ce{CH3COO-} ($\lambda=4{,}09$), les ions \ce{Na+}
étant spectateurs : la conductivité \res{\text{diminue}}.
\item Vers \res{\qty{530}{nm}}, au maximum d'absorption de l'ion permanganate : la mesure est la plus sensible et les
autres espèces, incolores, n'absorbent pas.
\item Mesures en continu et nombreuses, sans prélèvement ni perturbation du système, acquisition automatisable.
\end{questions}
\end{corrige}
\end{blocLivre}
\begin{blocLivre}{cor-4-8}
\begin{corrige}{8}{Suivi par titrages successifs}
\demarche{Écrire la relation à l'équivalence, puis tester la constance du quotient des concentrations sur des durées égales.}
\begin{questions}
\item La dilution et le refroidissement ralentissent fortement la décomposition : c'est une \res{\text{trempe}}, qui fige
la composition du prélèvement pendant le titrage. Les ions fer(III) sont un \res{\text{catalyseur}} (catalyse homogène).
\item À l'équivalence, $\frac{n(\ce{H2O2})}{5}=\frac{c\,V_\mathrm{E}}{2}$, d'où
$[\ce{H2O2}]=\frac{5\,c\,V_\mathrm{E}}{2\times\qty{10,0}{mL}}=\res{\num{5.00e-3}\times V_\mathrm{E}}$. En
\unit{mol.L^{-1}} : 0,0800 ; 0,0600 ; 0,0450 ; 0,0335 ; 0,0255 ; 0,0140 ; 0,0080 ; 0,0025.
\item $[\ce{H2O2}]=\qty{0,0400}{mol.L^{-1}}$ entre 10 et 15 min : $\res{t_{1/2}\approx\qty{12}{min}}$.
\item $\frac{0{,}0600}{0{,}0800}=0{,}75$ ; $\frac{0{,}0450}{0{,}0600}=0{,}75$ ; $\frac{0{,}0335}{0{,}0450}\approx0{,}74$.
Sur des durées égales, la concentration est multipliée par le même facteur : elle décroît exponentiellement, ce qui
caractérise une \res{\text{loi d'ordre 1}}. Avec $\mathrm{e}^{-5k}=0{,}75$ :
$k=\frac{-\ln0{,}75}{5}\approx\res{\qty{0,058}{min^{-1}}}$, et $\frac{\ln2}{k}\approx\qty{12}{min}$, en accord avec la question 3.
\end{questions}
\schema[corrige]{04-ex08-cor}
\end{corrige}
\end{blocLivre}
\begin{blocLivre}{cor-4-9}
\begin{corrige}{9}{La caféine dans l'organisme}
\demarche{Linéariser : si $C=C_0\,\mathrm{e}^{-kt}$, alors $\ln C$ est une fonction affine de $t$, de coefficient directeur $-k$.}
\begin{questions}
\item $\ln C$ : 1,275 ; 1,144 ; 0,867 ; 0,593 ; 0,315 ; $-0{,}261$.
\item Les points sont alignés sur une droite : $\ln C$ est une fonction affine décroissante du temps, donc $C$ décroît
exponentiellement : \res{\text{loi d'ordre 1}}. La modélisation donne $\ln C=1{,}42-0{,}140\,t$, d'où
$\res{k\approx\qty{0,14}{h^{-1}}}$.
\item $t_{1/2}=\frac{\ln2}{k}\approx\res{\qty{5,0}{h}}$ : valeur comprise entre 2 et 8 h, un peu supérieure à la moyenne
de 4 h : \res{\text{c'est cohérent}}.
\item 23 h correspond à $t=\qty{6}{h}$, soit 5 h après le maximum, environ une demi-vie : il reste
\res{\text{la moitié}} de la concentration maximale (\qty{1,81}{mg.L^{-1}} contre \qty{3,58}{mg.L^{-1}}). Au coucher,
la caféine encore présente peut retarder l'endormissement.
\end{questions}
\schema[corrige]{04-ex09-cor}
\end{corrige}
\end{blocLivre}
\begin{blocLivre}{cor-4-10}
\begin{corrige}{10}{Une transformation qui s'accélère}
\demarche{Suivre l'évolution de la pente de la courbe, puis confronter cette évolution aux facteurs cinétiques connus.}
\begin{questions}
\item La pente, d'abord faible, devient maximale vers $t=\qty{6}{min}$ (point d'inflexion), puis tend vers 0. La
tangente en $t=\qty{6}{min}$ passe par les points $(\qty{3,6}{min}\,;\,2{,}0)$ et $(\qty{8,4}{min}\,;\,0)$, donc
\[v_{\max}=\frac{\qty{2,0}{mmol.L^{-1}}}{\qty{4,8}{min}}\approx\res{\qty{0,42}{mmol.L^{-1}.min^{-1}}}.\]
\item \res{\text{Non}} : les concentrations des réactifs diminuent dès le début ; à elles seules, elles feraient
diminuer la vitesse dès le début, alors qu'ici elle commence par augmenter.
\item Au début, il y a très peu d'ions \ce{Mn^2+} : la réaction est lente. Les ions \ce{Mn^2+} formés catalysent la
réaction, qui accélère ; en fin de réaction, les ions permanganate s'épuisent et la vitesse diminue.
\item $\res{\dfrac{v_\mathrm{disp}(\ce{MnO4-})}{2}=\dfrac{v_\mathrm{disp}(\ce{H2C2O4})}{5}=\dfrac{v_\mathrm{app}(\ce{CO2})}{10}}$.
\item La vitesse serait maximale dès le départ : la courbe décroîtrait \res{\text{sans palier initial}}, comme pour une
transformation ordinaire.
\end{questions}
\schema[corrige]{04-ex10}
\end{corrige}
\end{blocLivre}
\begin{blocLivre}{cor-4-11}
\begin{corrige}{11}{Programmer un test de loi de vitesse}
\demarche{Une loi d'ordre 1 se traduit par un quotient $v/[\ce{CV+}]$ constant, égal à $k$.}
\begin{questions}
\item \code{v} est le taux de variation de la concentration entre les dates $t_{i-1}$ et $t_{i+1}$, changé de signe :
il approche la \res{\text{vitesse volumique de disparition}} à la date $t_i$. Le premier et le dernier point n'ont pas
de voisin d'un côté.
\item Les quotients $v/[\ce{CV+}]$ varient entre 0,0234 et 0,0259 \unit{min^{-1}} : ils sont pratiquement constants
(les écarts, de quelques pour cent, s'expliquent par les incertitudes de mesure). Donc $v=k\,[\ce{CV+}]$ :
\res{\text{loi d'ordre 1}}, avec $\res{k\approx\qty{0,025}{min^{-1}}}$ (moyenne des quotients).
\item $t_{1/2}=\frac{\ln2}{k}\approx\res{\qty{28}{min}}$. À $t=\qty{28}{min}$, on mesure \qty{9,94}{\micro\mole\per\litre},
environ la moitié de \qty{20,00}{\micro\mole\per\litre} : c'est vérifié.
\item \code{import matplotlib.pyplot as plt}, puis \code{plt.plot(C, V, "o")} et \code{plt.show()} : on attend des
points alignés sur une droite passant par l'origine.
\end{questions}
\schema[corrige]{04-ex11-cor}
\end{corrige}
\end{blocLivre}
\begin{blocLivre}{cor-4-12}
\begin{corrige}{12}{Décolorer un colorant alimentaire}
\demarche{Lire les dates où l'absorbance est divisée par deux, puis tester la constance des temps de demi-réaction successifs.}
\begin{questions}
\item Le colorant est la seule espèce colorée et, d'après la loi de Beer-Lambert, l'absorbance est proportionnelle à sa
concentration. On travaille à $\lambda_{\max}=\qty{630}{nm}$, où la mesure est la plus sensible.
\item $A=0{,}400$, moitié de 0,800, est atteinte pour $\res{t_{1/2}=\qty{30}{min}}$ (expérience 1) et
$\res{t_{1/2}=\qty{15}{min}}$ (expérience 2). Doubler la concentration en ions hypochlorite divise $t_{1/2}$ par deux :
la \res{\text{concentration d'un réactif}} est un facteur cinétique.
\item Dans l'expérience 1, $A$ passe de 0,800 à 0,400 en 30 min, puis de 0,400 à 0,200 en 30 min encore : les temps
de demi-réaction successifs sont égaux, \res{\text{la loi est d'ordre 1}}. Alors
$k_1=\frac{\ln2}{30}\approx\res{\qty{2,3e-2}{min^{-1}}}$.
\item $k_2=\frac{\ln2}{15}=2\,k_1$ : la constante double quand $[\ce{ClO-}]$ double, \res{\text{c'est cohérent}}. Pour
\qty{2,0e-2}{mol.L^{-1}}, $k=k_1/2$ et $\res{t_{1/2}=\qty{60}{min}}$.
\end{questions}
\schema[corrige]{04-ex12}
\end{corrige}
\end{blocLivre}
\begin{blocLivre}{cor-4-13}
\begin{corrige}{13}{Catalase, ions fer(III) ou rien ?}
\demarche{Comparer chaque expérience au témoin, puis exploiter le volume final de gaz, identique pour tous les catalyseurs.}
\begin{questions}
\item C'est un \res{\text{témoin}} : il montre que, sans catalyseur, la décomposition est très lente, dans les mêmes
conditions de volume et de température.
\item B : les ions fer(III) sont dissous, \res{\text{catalyse homogène}} ; C : la catalase est une enzyme,
\res{\text{catalyse enzymatique}}. La catalase est \res{\text{la plus efficace}} : le volume final est atteint en
quelques minutes, contre plus de 30 min avec les ions fer(III).
\item L'ébullition dénature l'enzyme, une protéine : sa structure est détruite et elle perd son activité catalytique.
L'expérience D se comporte comme le témoin.
\item $V_\mathrm{f}=\qty{48}{mL}$, donc $n(\ce{O2})=\dfrac{V_\mathrm{f}}{V_\mathrm{m}}=\qty{2,0e-3}{mol}$ et
$n(\ce{H2O2})_0=2\,n(\ce{O2})=\qty{4,0e-3}{mol}$, d'où
\[[\ce{H2O2}]_0=\frac{\qty{4,0e-3}{mol}}{\qty{20,0e-3}{L}}=\res{\qty{0,20}{mol.L^{-1}}}.\]
\item \res{\qty{48}{mL}} : un catalyseur ne modifie pas l'état final ; sans catalyseur, il faut seulement beaucoup plus de temps.
\end{questions}
\end{corrige}
\end{blocLivre}
\begin{blocLivre}{cor-4-14}
\begin{corrige}{14}{Décontaminer des lentilles de contact}
\demarche{Compter le nombre de demi-vies nécessaires pour passer de 3 \% à 0,01 \%, puis le convertir en durée.}
\begin{questions}
\item Le platine \res{\text{catalyse}} la décomposition du peroxyde d'hydrogène ; solide au contact d'une solution, il
réalise une \res{\text{catalyse hétérogène}}. L'évent laisse s'échapper le dioxygène formé et évite une surpression.
\item La teneur doit être divisée par $\frac{3}{0{,}01}=300=2^n$, soit $n=\frac{\ln300}{\ln2}\approx8{,}2$ temps de
demi-réaction : $8{,}2\times\qty{30}{min}\approx\res{\qty{4,1}{h}}$. Les 6 h imposées sont \res{\text{suffisantes}},
avec une marge de sécurité.
\item $6~\text{h}=12\times\qty{30}{min}$ : la teneur vaut $\frac{3~\%}{2^{12}}\approx\res{\num{7.3e-4}~\%}$.
\item Il faut alors $8{,}2\times\qty{45}{min}\approx\qty{6,2}{h}$ : au bout de 6 h, la teneur ($\frac{3~\%}{2^{8}}\approx0{,}012~\%$)
dépasse encore 0,01~\%, \res{\text{les lentilles irriteraient l'œil}}. Changer d'étui garantit un disque actif.
\end{questions}
\end{corrige}
\end{blocLivre}
\begin{blocLivre}{cor-4-15}
\begin{corrige}{15}{Hydrolyse du chlorure de tertiobutyle}
\demarche{Relier la conductivité à l'avancement par la loi de Kohlrausch, puis linéariser la loi d'ordre 1.}
\begin{questions}
\item La réaction forme des ions (\ce{H3O+}, \ce{Cl-}) à partir de molécules : la conductivité augmente avec l'avancement.
\item À la date $t$, $[\ce{H3O+}]=[\ce{Cl-}]=x/V$, donc $\sigma=\left(\lambda_{\ce{H3O+}}+\lambda_{\ce{Cl-}}\right)\frac{x}{V}$.
La transformation étant totale, $\sigma_\infty=\left(\lambda_{\ce{H3O+}}+\lambda_{\ce{Cl-}}\right)c_0$. Ainsi
$\frac{x}{V}=c_0\,\frac{\sigma}{\sigma_\infty}$ et $[\mathrm{RCl}]=c_0-\frac{x}{V}=\res{c_0\left(1-\sigma/\sigma_\infty\right)}$.
\item \code{y} vaut $\ln([\mathrm{RCl}]/c_0)$. Le programme le modélise par une fonction affine de coefficient directeur
$a=\qty{-2,328e-3}{s^{-1}}$ et d'ordonnée à l'origine $b=0{,}006\approx0$, en accord avec $\ln(1)=0$ :
$\ln([\mathrm{RCl}]/c_0)=-k\,t$, \res{\text{loi d'ordre 1}}, avec $\res{k\approx\qty{2,33e-3}{s^{-1}}}$ (les points
sont bien alignés, voir le graphique).
\item $t_{1/2}=\frac{\ln2}{k}\approx\res{\qty{3,0e2}{s}}$. À cette date, $\sigma=\sigma_\infty/2=\qty{0,81}{mS.cm^{-1}}$ :
c'est bien la valeur lue sur la courbe vers 300 s.
\item $k'=\frac{\ln2}{120}\approx\qty{5,8e-3}{s^{-1}}$, soit $\frac{k'}{k}=\frac{300}{120}=\res{2{,}5}$ : une hausse de
10 °C multiplie la vitesse par 2,5 ; la \res{\text{température}} est un facteur cinétique.
\end{questions}
\schema[corrige]{04-ex15}
\schema[corrige]{04-ex15-cor}
\end{corrige}
\end{blocLivre}
